\chapter{The menus}\label{cap:menus}
\subsection{The main menu}
\pantalla[0.62]{menu_principal_en}{The main menu in version \versionjuego.}
With \pant{ENGLISH} selected at start-up, the main menu has four options:
\begin{description}
\item[\pant{PLAY}] Start an adventure. Choose this last, once the characters are ready.
\item[\pant{BUY EQUIPMENT}] Visit the shop, where heroes spend gold on weapons and armour.
\item[\pant{CREATE HERO}] Decide which heroes take part and give them names.
\item[\pant{OPTIONS}] Music and typeface settings.
\end{description}
\subsection{Creating characters}
\pantalla[0.62]{crear_personaje_en}{The character screen: each hero's class on the left and name on the right.}
The \pant{CHOOSE HERO} screen shows the four heroes. To each one's right is their name---initially Sigmar the barbarian, Grugni the dwarf, Ladril the elf and Zoltan the wizard---or \pant{DEAD}. A dead character does not take part in the next adventure.
\begin{itemize}
\item Use \textbf{up} and \textbf{down}, or point with the mouse, to highlight a hero.
\item Press \textbf{right} to switch between the name and \pant{DEAD}. Use this to remove a hero who will not be playing, or to create a new character in a fallen hero's place. The new character starts with full points but no gold or equipment.
\item Press action on an active hero to open their menu: \pant{HERO NAME}, \pant{LOAD} and \pant{SAVE}.
\item \pant{BACK} returns to the main menu.
\end{itemize}
\pant{HERO NAME} shows an alphabet at the bottom of the screen. Choose up to six letters, one at a time, and finish with \pant{END}.

In the original game, \pant{LOAD} and \pant{SAVE} loaded and saved characters on tape. In this cartridge they use \textbf{codes}, explained below.

If every hero is dead, the game warns you and asks you to create new characters before playing again.
\subsection{Hero codes}\label{sec:codigos}
The cartridge cannot save data, so each hero is preserved with a 16-letter code: 14 data letters and two check letters. Old 14-letter codes are no longer accepted. The code contains:
\begin{itemize}
\item the hero's six-letter name;
\item gold;
\item current body points;
\item purchased weapons and armour, and the tool kit;
\item quest treasures found during adventures.
\end{itemize}
It does not save potions or spells, which are lost at the end of each adventure, or the weapon currently held: after restoring a hero, select it again in the inventory.
\dospantallas{codigo_guardar_en}{codigo_cargar_en}{Left: \pant{SAVE} shows the hero's code. Right: \pant{LOAD} asks for it letter by letter.}
\subsubsection*{Saving: display the code}
In \pant{CREATE HERO}, select a hero, press action and choose \pant{SAVE}. The screen shows the code in three letter groups, for example \texttt{QGYNT ZDWQL PYZGTI}. Copy it exactly and press action to return. Do this for all four heroes after an adventure and after visiting the shop: the code describes the hero at the moment you request it.
\subsubsection*{Loading: enter the code}
Select a hero of the same class---a barbarian's code only works for the barbarian---press action and choose \pant{LOAD}. Enter letters in any of three ways:
\begin{itemize}
\item type them directly on the keyboard;
\item point at each letter in the alphabet and click the mouse button;
\item move along the alphabet with \textbf{left} and \textbf{right}, then press action.
\end{itemize}
\pant{END} deletes the last letter; when the line is empty, it returns to the character screen without changing anything. Once letter 16 is entered, the hero appears with their name, gold and equipment. If a letter is wrong, the game displays \pant{INVALID CODE} and lets you try again. A mistyped code is very unlikely to be accepted.
\aviso{Warning}{Loading a code replaces the hero already in that slot. To preserve that hero, write down their code first.}
\subsection{Buying equipment}
\pantalla[0.7]{tienda_en}{The shop. Buyers are listed above, weapons and armour with their prices below.}
During adventures, heroes find gold and receive rewards. They can spend it in the shop on better equipment, whose benefits apply automatically.
\begin{enumerate}
\item Select the buyer from the upper list. Their gold is displayed.
\item Click the item you want. If the hero has enough gold and their class can use it, the item becomes theirs.
\item \pant{BACK} in the item list ends that hero's shopping; \pant{BACK} in the hero list returns to the main menu.
\end{enumerate}
\begin{table}[H]\centering\small
\begin{tabularx}{\linewidth}{@{}lrL@{}}
\toprule\textbf{Weapon (screen label)} & \textbf{Price} & \textbf{Effect}\\\midrule
Staff (\pant{STAFF}) &100&2 attack dice; may also attack diagonally\\
Dagger (\pant{DAGGER}) &150&2 attack dice; may also attack diagonally\\
Broadsword (\pant{BROADSWORD}) &250&3 attack dice\\
Battle axe (\pant{BATTLE AXE}) &400&4 attack dice; prevents use of a shield\\
Crossbow (\pant{CROSSBOW}) &350&3 ranged attack dice; cannot target an adjacent enemy\\
Hand axe (\pant{HAND AXE}) &150&2 attack dice; may be thrown, but is then lost\\
Spear (\pant{SPEAR}) &150&2 attack dice; may attack diagonally or be thrown, but is lost when thrown\\\bottomrule
\end{tabularx}\caption{Weapons for sale.}\end{table}
\begin{table}[H]\centering\small
\begin{tabularx}{\linewidth}{@{}lrL@{}}
\toprule\textbf{Armour/equipment} & \textbf{Price} & \textbf{Effect}\\\midrule
Shield (\pant{SHIELD}) &100&One additional defence die\\
Helmet (\pant{HELM}) &120&One additional defence die\\
Chain mail (\pant{CHAIN MAIL}) &450&3 defence dice\\
Plate armour (\pant{PLATE}) &850&4 defence dice, but its weight reduces movement to one die\\
Tool kit (\pant{KIT}) &250&Disarm discovered traps; the dwarf already has one\\\bottomrule
\end{tabularx}\caption{Armour and equipment for sale.}\end{table}
The \textbf{wizard} can buy only the staff and tool kit: the other weapons and armour are not for him.
\subsection{Options}
\pantalla[0.62]{controles_en}{The options screen.}
Operate this like the main menu: \textbf{up} and \textbf{down} move the bar and action selects, or point at an option with the mouse and click.
\begin{description}
\item[\pant{MUSIC ON}/\pant{MUSIC OFF}] Turn music on or off.
\item[\pant{GOTHIC FONT}/\pant{NORMAL FONT}] Switch the whole game's typeface between the default Gothic lettering and a plainer, easier-to-read font.
\item[\pant{BACK}] Return to the main menu.
\end{description}
\subsection{Playing: spells and adventure selection}
After choosing \pant{PLAY}, the magic-using characters share out the spells before you choose an adventure. There are four groups---Air, Fire, Earth and Water (\pant{AIR SPELLS}, \pant{FIRE SPELLS}, \pant{EARTH SPELLS}, \pant{WATER SPELLS})---each containing three spells.
\pantalla[0.7]{hechizos_en}{Choosing the spell groups.}
The heading says whose choice it is: \pant{WIZARD CHOOSE SPELLS} or \pant{ELF CHOOSE SPELLS}.

The wizard chooses one group first. The elf chooses next, and the two remaining groups go to the wizard. If only one of them takes part, that hero chooses alone: one group for the elf, three for the wizard.

The adventure list then appears, split across two pages:
\dospantallas{aventuras1_en}{aventuras2_en}{The two adventure pages: the original game and \textit{Return of the Witch Lord}.}
\begin{itemize}
\item The first contains the original game's \textbf{fourteen adventures}. Its last line, \pant{NEXT}, opens the second page.
\item The second contains the expansion's \textbf{ten adventures}. Its last line, \pant{BACK}, returns to the first.
\end{itemize}
With a mouse, the mission you point at is highlighted directly.

Choose an adventure to read Mentor's briefing. Read it carefully: it explains your objective. Press action to enter the dungeon. Play the adventures in order: each is slightly harder than the last and the story progresses between them.
\pantalla[0.7]{pergamino1_en}{Mentor introduces the first adventure.}
